\exercises

\tierunderstand

\begin{bt}
  \item Phương trình~\eqref{eq:ch05-phan-ra} không có \(\vect{x}\) ở vế phải.
    Viết ra chính xác chỗ nào trong mô hình làm cho chuyện đó đúng, rồi trả
    lời: nếu decoder được đọc lại \(\vect{v}\) ở mọi bước thay vì chỉ ở bước
    khởi tạo, vế phải đổi thế nào và \tn{chỗ thắt}{bottleneck} dịch đi đâu?

  \item Bài viết ô nhớ của họ mang 8000 số thực cho một câu. Dẫn ra con số ấy
    từ cấu hình mà mục 3.4 mô tả. Rồi tính con số tương ứng cho mô hình ở
    dòng cuối bảng của mục~\ref{sec:ch05-bop-vector}, và nói tỷ số giữa hai
    con số là bao nhiêu.

  \item Trong bảng ở mục~\ref{sec:ch05-bop-vector}, dòng \(d_h = 64\) cho câu
    ngắn 0.9797 và câu dài 0.0724. Giải thích chênh lệch ấy chỉ bằng bề rộng
    \tn{vector ngữ cảnh}{context vector}, rồi nói xem lời giải thích của bạn dự
    đoán gì cho một corpus mà câu dài gấp đôi corpus này.

  \item Lập luận của bài về phép đảo nói khoảng cách trung bình giữa các cặp từ
    tương ứng \emph{không đổi}. Chứng minh câu đó cho một câu \(n\) từ với
    đồng chỉnh một-một theo thứ tự. Rồi chỉ ra đại lượng nào thật sự đổi, và
    vì sao đại lượng đó mới là đại lượng chương~\ref{ch:lstm} quan tâm.

  \item Mục~\ref{sec:ch05-giai-ma} nói beam rộng hơn không bảo đảm trả về
    \tn{giả thuyết}{hypothesis} có log xác suất cao hơn. Dựng một phản ví dụ nhỏ nhất bạn nghĩ ra
    được: một \tn{từ điển}{vocabulary} ba token, một cây xác suất bạn tự đặt,
    và hai bề rộng
    beam cho hai kết quả khác nhau theo chiều đó.

  \item Ensemble lấy trung bình log xác suất của các mô hình ở mỗi bước, chứ
    không bỏ phiếu chọn câu cuối. Nêu một tình huống cụ thể mà hai cách cho hai
    câu khác nhau, rồi nói cách nào khớp với \eqref{eq:ch05-argmax} hơn và vì
    sao.

  \item Corpus ở mục~\ref{sec:ch05-corpus} có \tn{câu đích}{target sentence}
    dài hơn \tn{câu nguồn}{source sentence} ở 0.7800 số cặp. Kể tên hai cơ chế
    trong văn phạm gây ra chuyện đó, rồi nói cơ chế nào trong hai cái đó cũng
    có trong cặp Anh-Pháp mà bài dùng.
\end{bt}

\tierapply

Trạng thái xuất phát cho cả năm bài: clone repo companion, đứng ở thư mục gốc
của nó, và checkout đúng tag của chương này.

\begin{minted}{console}
$ git clone https://github.com/Giang-Dang/rnn-to-transformer-lab.git
$ cd rnn-to-transformer-lab
$ git checkout ch05
$ conda env create -f environment.yml
$ conda activate rnn-to-transformer-lab
$ python verify.py --only ch05
\end{minted}

Dòng cuối phải in ra \code{verify: ok}. Nếu không, đừng làm tiếp: mọi con số
bạn đo sau đó sẽ không so được với sách. Nếu output hiện ra toàn dấu hỏi thay
cho chữ tiếng Việt thì console của bạn không ở UTF-8; số vẫn đúng, chỉ chữ
hỏng.

\begin{bt}
  \item Bảng ở mục~\ref{sec:ch05-bop-vector} dừng ở \(d_h = 128\). Thêm 256 vào
    \code{WIDTHS} trong \code{ch05\_bottleneck.py} rồi chạy lại. Cột
    \code{long} đi tới đâu, và thời gian chạy đội lên bao nhiêu? Ngân sách của
    script là 360 giây; nói rõ bạn đã bỏ dòng nào đi để đo tiếp nếu phải bỏ.

  \item \code{encode} đóng băng trạng thái ở vị trí đệm bằng hai dòng
    \code{live}. Bỏ chúng đi, cho \code{c\_next, h\_next} chạy thẳng, rồi chạy
    \code{python -m pytest -q tests/test\_ch05.py} và nói test nào hỏng. Sau đó
    chạy lại \code{ch05\_bottleneck.py} ở riêng \(d_h = 128\) và nói độ chính
    xác đổi bao nhiêu. Con số ấy nhỏ hay lớn hơn bạn đoán?

  \item Mục~\ref{sec:ch05-giai-ma} kể một luật dừng sai làm beam 2 kém giải mã
    tham lam. Dựng lại nó: trong \code{beam\_decode}, bỏ dòng cắt tỉa theo
    \code{best\_finished}, cho vòng lặp nạp lại \(\mathcal{L}\) đủ \(B\) phần
    tử bất kể bao nhiêu giả thuyết vừa kết thúc, rồi thoát khi
    \code{len(finished) >= beam}. Chạy lại bảng beam trên mô hình seed 0 và xác
    nhận 0.8400 ở beam 2, so với 0.8533 của tham lam. Rồi giải thích, bằng một
    câu test khẳng định được, vì sao chỗ mất nằm ở câu \emph{ngắn}.

  \item \code{ch05\_search.py} huấn luyện ba mô hình khác nhau ở seed. Đổi
    \code{SEEDS} thành năm seed, và nhớ đổi cả dãy \code{(1, 2, 3)} ở vòng lặp
    in bảng ensemble, vì nó được viết cứng và sẽ vẫn dừng ở ba. Ensemble năm có
    hơn ensemble ba không, và phần hơn ấy so với phần hơn khi đi từ một lên ba
    thì thế nào? Ghi lại thời gian chạy.

  \item Hiện corpus chỉ sinh câu tối đa hai mệnh đề. Trong
    \code{ch05\_bottleneck.py}, tìm lời gọi \code{disjoint\_splits} rồi đặt
    tham số \code{max\_clauses} của nó bằng 3. Giữ nguyên mọi thứ khác và chạy
    lại cả bảng. Ngưỡng \(d_h\) mà câu dài bắt đầu giải được dịch đi đâu, và
    chuyện đó ủng hộ hay bác bỏ cách đọc bảng ở
    mục~\ref{sec:ch05-bop-vector}?
\end{bt}

\tierextend

\begin{bt}
  \item Mục 5 của bài viết rằng họ không huấn luyện nổi một mạng hồi quy thường
    trên bài toán không đảo, nhưng tin rằng nó sẽ huấn luyện được khi câu nguồn
    đảo, và nói rõ họ không kiểm bằng thí nghiệm. Kiểm hộ họ trên corpus này:
    thay \code{LstmLayer} bằng một tầng hồi quy tanh thường trong cả encoder
    lẫn decoder, chạy cả hai vế đảo và không đảo, ba seed mỗi vế. Kết quả của
    bạn ủng hộ hay bác bỏ niềm tin ấy, và ở quy mô này thì câu trả lời có ý
    nghĩa gì đối với quy mô của họ?

  \item Bảng ở mục~\ref{sec:ch05-giai-ma} cho ba mô hình cùng thiết lập tản gần
    hai mươi điểm. Tìm xem cái gì trong lần huấn luyện của seed 1 khác hai seed
    kia: vẽ đường loss của cả ba, tìm chỗ chúng tách nhau, và xem có can thiệp
    nào rẻ hơn việc huấn luyện thêm hai mô hình mà kéo được seed 1 lên không.
    Nếu có thì ensemble còn đáng giá như bảng nói không?

  \item Mục~\ref{sec:ch05-con-lai-gi} phát biểu rằng chỗ thắt là một tỷ lệ giữa
    bề rộng vector và lượng thông tin phải đi qua, chứ không phải một tính chất
    của kiến trúc. Biến câu đó thành một phát biểu đo được: định nghĩa
    \enquote{lượng thông tin của một câu} trên corpus này bằng một đại lượng
    tính được, rồi kiểm xem ngưỡng \(d_h\) mà độ chính xác vượt một mức cho
    trước có tỷ lệ với đại lượng ấy không. Nói rõ định nghĩa của bạn hỏng ở đâu.

  \item Đây là chỗ còn mở thật sự. \(\vect{v}\) là \(2 d_h\) số dấu phẩy động,
    nên về lý thuyết nó chứa được rất nhiều bit, thừa sức mã hóa mọi câu của
    corpus này ngay ở \(d_h = 8\). Vậy mà \(d_h = 8\) cho 0.0000. Nghĩa là giới
    hạn không nằm ở dung lượng biểu diễn mà ở chỗ khác: ở cái mà huấn luyện
    bằng gradient tìm được, hoặc ở chỗ decoder giải mã được. Tách hai khả năng
    ấy ra bằng một thí nghiệm, và biết trước rằng chưa có phát biểu nào được
    chấp nhận rộng rãi về việc một vector trạng thái học được thật sự giữ được
    bao nhiêu.
\end{bt}
